\section{Discussion}\label{sec:discussion}

\subsection{Why the ten surveys could not separate the method from the line}

The held-out comparison of~\cite{cb002} found the method level with the line in the setting on power. The theory and the study
say why, term by term. The engine's throughput is affine, so on throughput the method is a ridge regression and its only gain
is the shrinkage of the plane, which is real at ten surveys but is a variance saving that a regression in both regressors
also mostly obtains. On power the engine's curvature over the range of the surveys is \eRmsS~kW against a noise of
\sigKw~kW, so even when it is perfectly aligned it removes \curvShareKw\pct{} of the noise variance. The feed size explains
little of the power beyond the setting ($\phi^2=\phiTwoKwS$~kW$^2$), so the third regressor of the discrepancy costs about as
much variance as it earns, and the crossing size $n^\star=\nStarKwS$ lies just above ten. On a virtual plant fitted to the ten
surveys, with the true scalars handed to it, the method ties the line at ten surveys ($\lineOrTen$~kW$^2$ in paired squared
error) and overtakes it from twenty on; with learned scalars from forty. A real campaign of ten surveys could not have shown a
difference that the theory says is a few kW$^2$ against a noise variance of about \sigKw$^2$.

\subsection{When the physics pays}

Corollary~\ref{cor:align} reduces the question to two numbers over the operating range of interest: the alignment $\rho$ of the
engine's departure from a plane with the plant's, and their relative size $\kappa$. Both can be estimated before a campaign.
The engine's departure $e$ is computable at any operating point. The plant's departure $r$ is not, but its size has a lower
bound in the scatter of past surveys about a plane, which contains the noise, and a sign that process knowledge often gives:
whether the capacity is concave or convex in the setting. The capacity of the engine used here is affine in the setting at the
fixed speed, so the physics in it cannot help predict throughput; a capacity law curved like the plant's, such as the power
law of Rocha et al.~\cite{rocha} or the concave capacity that Duarte et al.~\cite{duarte} measured on an HP400 over 6~mm of
setting, is the condition for any persistent gain on throughput. When the condition holds the gain persists as surveys
accumulate and grows under extrapolation, which is where the physics is meant to be used; its size is bounded by the plant's
own curvature relative to the noise, \curvCeilTph\pct{} of the squared error for the curvature the published capacity law gives
over the surveys window.

The second lesson is that the discrepancy prior is a predictive decision. Theorem~\ref{thm:fixed} gives the risk-optimal scale
in closed form, and on the plant fitted to the ten surveys the product's scale is close to it. The orthogonal prior, introduced
to give the scalars a meaning, takes the room of the discrepancy away when the engine has more sensitivities than the
discrepancy has regressors, and its predictions degrade toward those of the engine without a discrepancy. A practitioner who
wants interpretable scalars and good predictions can calibrate twice, with the orthogonal prior for the scalars and with the
standard prior for the predictions, as Proposition~\ref{prop:priors}(a) shows both are the same linear model with different
penalties.

\subsection{What to measure next}

The design study answers the question the earlier paper left open, which measurement would teach the scalars most. A product
size measured with a noise of 2 to 5\pct{} is worth as much as doubling the campaign or more, and it is the only added output
that reaches the classification and breakage scalars, which throughput and power leave near their priors. More surveys at the
same speed cannot narrow the capacity base below the floor of~\eqref{eq:floor}, which the ten surveys already approach
(\ceilShare\pct{} of the possible narrowing), and varying the speed by ten percent around the engine's capacity optimum adds
nothing, because the capacity is flat there. A design that moved the capacity base would need a factor that changes the
capacity outside the span of the discrepancy regressors by an amount comparable to the noise, which the speed does not do on
this engine; the study did not evaluate other factors. For prediction rather than identification, the cheapest improvement is
the noise itself: every gain of the method is measured against $\sigma^2$, so a better estimate of the throughput than the feed
rate of a survey, or of the power than the one estimated from the motor current, raises both the curvature and the plane gain in
proportion.

\subsection{Negative results}

Three hypotheses failed in part. The speed clause of H5 failed for the reason just given. H4 failed for the orthogonal prior,
which changes the risk more than any other component except the discrepancy itself. And the engine with half the exponents of
the plant's capacity law, a curvature of the right sign, lost to least squares on the wide window: the sign of a curvature does
not decide the gain, its projection on the plant's curvature over the prediction range does. Each of these is a statement about
this engine and these plants; each follows from a result of Section~\ref{sec:theory} that holds generally. The study also
confirms the negative result of~\cite{cb002} rather than overturning it: on a virtual plant built from the fit to the ten
surveys, with the scalars handed to it and the noise known, the method still ties the line at ten surveys. No change of prior,
scale or sampler would have changed the earlier comparison; only more surveys, less noise or an engine with a curved capacity
would.
